% Generated draft: requires mathematical and source review.
\documentclass[11pt]{article}
\usepackage[margin=25mm]{geometry}
\usepackage{fontspec}
\setmainfont{Latin Modern Roman}
\usepackage{amsmath,amssymb,mathtools,unicode-math}
\setmathfont{Latin Modern Math}
\usepackage{enumitem,needspace,xurl,hyperref}
\setlength{\parindent}{0pt}
\setlength{\parskip}{5pt}
\setlength{\emergencystretch}{3em}
\newcommand{\N}{\mathbb{N}}
\newcommand{\Z}{\mathbb{Z}}
\newcommand{\Q}{\mathbb{Q}}
\newcommand{\R}{\mathbb{R}}
\newcommand{\C}{\mathbb{C}}
\newcommand{\F}{\mathbb{F}}
\newcommand{\A}{\mathbb{A}}
\newcommand{\PP}{\mathbb{P}}
\newcommand{\E}{\mathbb{E}}
\newcommand{\eps}{\varepsilon}
\newcommand{\dd}{\,\mathrm{d}}
\begin{document}
\section*{UnsolvedMath ID: 30001842}
\textbf{Source problem code:} catalogue entry
\subsection{Definitions and mathematical statement}
For a prime \(p\), a nonsplit Cartan subgroup \(C_{\mathrm{ns}}(p)\subseteq\operatorname{GL}_2(\F_p)\) is the image of \(\F_{p^2}^{\times}\) acting by multiplication on the two-dimensional \(\F_p\)-vector space \(\F_{p^2}\), after choosing a basis. Different choices of basis give conjugate subgroups. Under the usual convention, \(X_{\mathrm{ns}}(p)\) is the smooth projective modular curve over \(\Q\) associated with this subgroup, obtained by compactifying the moduli curve of elliptic curves with full level-\(p\) structure modulo \(C_{\mathrm{ns}}(p)\).

A noncuspidal geometric point of this curve is a \emph{CM point} if its underlying elliptic curve \(E/\overline{\Q}\) has complex multiplication, meaning that \(\operatorname{End}_{\overline{\Q}}(E)\) strictly contains \(\Z\), or equivalently is an order in an imaginary quadratic field. Define
\[
\operatorname{CM}_{\Q}(X_{\mathrm{ns}}(p))
=\{x\in X_{\mathrm{ns}}(p)(\Q):x\text{ is noncuspidal and its underlying elliptic curve has CM}\}.
\]
The question in the supplied record is whether
\[
\exists\,c\in\Z_{\geq 2}\quad
\forall\,p\text{ prime},\qquad
p>c\ \Longrightarrow\ 
X_{\mathrm{ns}}(p)(\Q)=\operatorname{CM}_{\Q}(X_{\mathrm{ns}}(p)).
\]
Taking an integer bound at least \(2\) does not change the existence question. The bound is uniform over all primes, and the assertion does not claim that any rational CM points exist.

\emph{Unresolved source qualifications.} The supplied record does not define its modular-curve notation. The definition above uses the nonsplit Cartan subgroup itself; replacing it by its normalizer gives the curve commonly denoted \(X_{\mathrm{ns}}^{+}(p)\), which is a different formulation. The intended convention requires confirmation against the source. Moreover, the statement as supplied says all rational points on \(X_{\mathrm{ns}}(p)\), without excluding cusps. On the projective convention above, the assertion therefore also requires that there be no rational cusps for primes exceeding the bound. If the intended target concerns only noncuspidal rational points, that restriction must be confirmed rather than silently inserted.
\subsection{Short English statement}
Is there a single bound beyond which every rational point on the indicated nonsplit Cartan modular curve is a noncuspidal point representing an elliptic curve with complex multiplication? The precise Cartan versus normalizer convention, and whether cusps are included, require source confirmation. This is a statement-only draft; the supplied open status is a dated curation claim.
\subsection{Sources}
\begin{itemize}
\item Local UnsolvedMath registry, supplied complete record with permanent ID \(30001842\), in the Oberwolfach Reports Open Problems collection; exact source problem code \texttt{catalogue entry}. The project registry locations are \texttt{data/unsolvedmath/problems.json} and the supplementary \texttt{catalog/problems.json}; no independent file inspection is asserted here.
\item Supplied report citation: \emph{Explicit Methods in Number Theory (2012)}, Oberwolfach Reports, DOI \texttt{10.4171/owr/2011/35}, pp.~1971--2044. Supplied source URL and literature-source link labelled ``original OWR report'': \url{https://doi.org/10.4171/owr/2011/35}. The record gives proposed year \(2011\) and supplies no named proposer.
\item Supplied literature-source link labelled ``level-13 rational-point classification'': \url{https://arxiv.org/abs/1711.05846}. Its precise applicability to the modular-curve convention in this question requires review.
\item Literature assessment dated exactly \texttt{2026-08-21}: the supplied curation labels the problem open and states that rational-point classifications are known for selected small levels, including level \(13\), while the asserted uniform statement for all sufficiently large primes is not proved. This is supplied dated curation, not an independently verified literature assessment. The background explicitly requests independent verification of time-sensitive status and cited claims.
\item The supplied statement-status label is \texttt{corrected\_verified}. Its accompanying verification note claims comparison with the official Oberwolfach report and normalization of notation, context, and question form without changing the source-level claim. This is provenance supplied by the registry, not verification performed for this draft.
\end{itemize}
\end{document}
